\chapter{Bối cảnh và Công trình liên quan}\label{chapter-2-background-and-related-work}

Chương này trình bày nền tảng bối cảnh mà khóa luận được xây dựng dựa trên đó: từ thiết lập học tăng cường (reinforcement learning) làm động lực cho self-distillation ra đời (§2.1), qua phương pháp SDPO và hàm mất mát của nó (§2.2), đến chế độ test-time cùng hệ đo lường (metric) tương ứng (§2.3), không gian thiết kế reprompt-template (§2.4), những cảnh báo về hiện tượng epistemic suppression (§2.5), một phương pháp self-distillation tương đồng (§2.6), và khép lại bằng khoảng trống nghiên cứu mà khóa luận này hướng đến giải quyết (§2.7).

\section{RLVR và GRPO}\label{rlvr-and-grpo}

Reinforcement learning with verifiable rewards (RLVR) hiện là phương pháp post-training chủ đạo cho các tác vụ lập trình (code) và toán: một bộ kiểm tra (checker) tự động cung cấp phần thưởng (reward), thay vì dựa vào human preference như trong reinforcement learning from human feedback {[}15{]} hay các bản đơn giản hóa của nó về tối ưu hóa sở thích (preference-optimization) {[}16{]}. Group Relative Policy Optimization (GRPO) {[}8{]} là một đại diện phổ biến nhất: phương pháp này lấy mẫu (sample) một nhóm các rollout cho mỗi bài toán, rồi gán cho mỗi rollout một advantage tương đối so với trung bình nhóm, \(A_{i,t} = r_i - \text{mean}_j\{r_j\}\), một hằng số được giữ nguyên cho mọi token trong một chuỗi.

Ưu điểm của GRPO là nó loại bỏ được value network riêng biệt mà các phương pháp kiểu PPO {[}14{]} vẫn cần: trung bình nhóm tự đóng vai trò là baseline, còn advantage được ước lượng trực tiếp từ độ phân tán của reward trong nhóm đã sample {[}8{]}. Nhờ vậy, phương pháp này được giữ ở mức tối giản, và đó cũng là một phần lý do nó trở thành lựa chọn mặc định cho các miền verifiable.

Tuy nhiên, chính sự tối giản đó lại là điểm yếu trong việc phân bổ tín hiệu học (credit assignment). Một scalar outcome reward chỉ cho biết liệu toàn bộ rollout có thành công hay không, chứ không ghi nhận được token nào thực sự chịu trách nhiệm, dẫn đến tín hiệu học rất thưa thớt {[}1{]}. Process supervision (chấm điểm từng bước lập luận trung gian thay vì chỉ đánh giá kết quả cuối cùng {[}22{]}) là một cách giải quyết sự thưa thớt này; trong khi đó, SDPO đi theo một hướng khác: trích xuất một target dày đặc (dense target) theo từng token ngay từ chính feedback. Điểm yếu này sẽ trở thành thất bại hoàn toàn khi mọi rollout trong một nhóm đều sai: advantage tương đối triệt tiêu về 0, không có gradient nào được truyền đi, và phương pháp không thể học trên bài toán đó cho đến khi nó tìm ra lời giải đúng ít nhất một lần. Ở những bài toán khó, nơi sự thành công vốn rất hiếm hoi, đây chính xác là rào cản khiến quá trình học bị đình trệ.

Một giải pháp khắc phục tự nhiên là distill từ một teacher mạnh hơn. Nhưng điều đó yêu cầu phải có sẵn một teacher vượt trội hơn student, trong khi mục tiêu ở đây là nâng cao giới hạn năng lực (capability ceiling) thay vì nén một mức năng lực đã biết {[}1{]}. Nghịch lý này chính là động lực để tìm kiếm một teacher không đến từ bên ngoài, mà được dẫn xuất từ chính bản thân student.

\section{SDPO}\label{sdpo}

\subsection{Self-teacher và rich feedback}\label{the-self-teacher-and-rich-feedback}

Hübotter và cộng sự {[}1{]} nhận ra rằng nhiều môi trường verifiable vốn đã cung cấp sẵn tokenized feedback: lỗi thời gian chạy (runtime errors), diff của test case thất bại, hay nhận xét của judge, nhưng những tín hiệu này lại bị lãng phí khi mọi thông tin bị thu gọn thành một scalar reward duy nhất. Họ đề xuất reinforcement learning with rich feedback (RLRF) như một khuôn khổ (framework) nhằm bảo toàn tín hiệu đó, và SDPO chính là một cách hiện thực hóa của khuôn khổ này. Cơ chế cốt lõi biến ý tưởng này thành hiện thực là in-context learning: cùng một mô hình, khi được điều kiện hóa thêm trên feedback \(f\) bên cạnh câu hỏi \(x\), thường có khả năng tự định vị được lỗi sai của chính mình. Mô hình được điều kiện hóa này, \(\pi_\theta(\cdot \mid x, f)\), đóng vai trò là \emph{self-teacher}; trong khi mô hình không điều kiện \(\pi_\theta(\cdot \mid x)\) đóng vai trò là student. Vì cả hai cùng chia sẻ chung tập trọng số \(\theta\), đây là quá trình self-distillation đúng nghĩa, kế thừa nền tảng của knowledge distillation {[}13{]}, bao gồm cả các biến thể mức chuỗi (sequence-level) cho ngôn ngữ {[}24{]}, chỉ khác ở chỗ nguồn lợi thế của teacher đến từ feedback thay vì từ một mạng neural lớn hơn.

\subsection{Hàm mất mát và credit assignment dày}\label{the-loss-and-dense-credit-assignment}

SDPO distill phân bố next-token mang tính hồi cứu (retrospective) của self-teacher vào student thông qua một hàm KL phân bổ theo từng token:

\[
\mathcal{L}_{\text{SDPO}}(\theta) = \sum_{t=1}^{|y|} \mathrm{KL}\,\Big(\pi_\theta(\cdot \mid x, y_{<t}) \,\big\|\, \mathrm{stopgrad}\,\pi_\theta(\cdot \mid x, f, y_{<t})\Big),
\]

trong đó \passthrough{\lstinline!stopgrad!} ngăn không cho teacher suy biến (collapse) về student và bỏ qua \(f\) {[}1{]}. Vì teacher được ngắt gradient (detached), đạo hàm tại một logit của student trở thành \(\partial\mathcal{L}/\partial z_v = \pi_S(v) - \pi_T(v)\), nghĩa là optimizer sẽ gia tăng giá trị của bất kỳ logit nào được teacher gán xác suất cao hơn so với student. So với GRPO, khác biệt lớn nhất nằm ở mật độ (density) của target:

\begin{longtable}[]{@{}
  >{\raggedright\arraybackslash}p{(\columnwidth - 4\tabcolsep) * \real{0.3333}}
  >{\raggedright\arraybackslash}p{(\columnwidth - 4\tabcolsep) * \real{0.3333}}
  >{\raggedright\arraybackslash}p{(\columnwidth - 4\tabcolsep) * \real{0.3333}}@{}}
\caption{Độ dày credit-assignment: target mỗi token của GRPO so với SDPO.}\\
\toprule\noalign{}
\begin{minipage}[b]{\linewidth}\raggedright
Phương pháp
\end{minipage} & \begin{minipage}[b]{\linewidth}\raggedright
Target mỗi token
\end{minipage} & \begin{minipage}[b]{\linewidth}\raggedright
Tín hiệu mỗi chuỗi
\end{minipage} \\
\midrule\noalign{}
\endfirsthead
\toprule\noalign{}
\begin{minipage}[b]{\linewidth}\raggedright
Phương pháp
\end{minipage} & \begin{minipage}[b]{\linewidth}\raggedright
Target mỗi token
\end{minipage} & \begin{minipage}[b]{\linewidth}\raggedright
Tín hiệu mỗi chuỗi
\end{minipage} \\
\midrule\noalign{}
\endhead
\bottomrule\noalign{}
\endlastfoot
GRPO & scalar advantage × one-hot & \(O(1)\) \\
SDPO (mức chuỗi) & scalar × soft distribution & \(O(1)\) \\
SDPO (mức token) & soft distribution, top token & \(O(T)\) \\
SDPO (mức logit) & soft distribution trên top-\(K\) vocab & \(O(T \cdot K)\) \\
\end{longtable}

\begin{figure}[H]
\centering
\begin{tikzpicture}[font=\footnotesize,>=Stealth,semithick,every node/.style={draw,rounded corners,align=center,minimum height=0.7cm,text width=3.3cm,inner sep=2pt}]
\node(gr) at (0,3){scalar reward $r$};
\node(ga) at (0,2){one advantage / sequence};
\node(gt) at (0,1){same signal every token ($O(1)$)};
\node(sf) at (4.4,3){feedback $f$};
\node(st) at (4.4,2){teacher per-token dist.};
\node(stk) at (4.4,1){soft top-$K$ target / position ($O(T\cdot K)$)};
\draw[->](gr)--(ga);\draw[->](ga)--(gt);\draw[->](sf)--(st);\draw[->](st)--(stk);
\node[draw=none,font=\bfseries] at (0,3.8){GRPO (sparse)};
\node[draw=none,font=\bfseries] at (4.4,3.8){SDPO (dense)};
\end{tikzpicture}
\caption{Độ dày credit-assignment, GRPO so với SDPO. GRPO gắn một scalar advantage cho cả chuỗi, nên mọi token nhận cùng tín hiệu (\(O(1)\)); SDPO điều kiện hóa teacher trên feedback và cấp một target mềm, \(K\)-chiều tại mỗi vị trí (\(O(T\cdot K)\)).}
\label{fig:credit-density}
\end{figure}

Ở hình thái đầy đủ nhất (mức logit), target là một phân bố \(K\)-chiều tại mọi vị trí, thay vì một scalar duy nhất cho toàn bộ chuỗi; đây chính là nguồn gốc tạo nên hiệu suất lấy mẫu (sample efficiency) vượt trội mà SDPO đạt được {[}1{]}. Để giới hạn chi phí tính toán trên một vocabulary khoảng 150k, phương pháp này chỉ giữ lại top-\(K\) logit của teacher (\(K{=}100\), hoặc 20 đối với cấu hình code {[}4{]}) cộng thêm một nhóm đuôi (tail bucket). Chiều (direction) của KL được tham số hóa qua \(\alpha\): forward KL (\(\alpha{=}0\)) có tính chất mode-covering, reverse KL (\(\alpha{=}1\)) có tính chất mode-seeking, còn JSD (\(\alpha{=}0.5\)) đóng vai trò nội suy giữa hai cực này; lựa chọn này tuân theo Agarwal và cộng sự {[}11{]} để đảm bảo tính ổn định. Self-teacher cũng cần được điều chuẩn (regularize), thường bằng một EMA của trọng số student hoặc một phép nội suy trust-region, vì nếu thiếu bước này, teacher sẽ phân kỳ (diverge) {[}1{]}. Cách hiện thực (implementation) dùng trong khóa luận này được đặc tả kỹ hơn ở §3.3.5; toàn bộ phần dẫn xuất toán học được đính kèm trong ghi chú dự án.

\subsection{Ba chế độ}\label{three-regimes}

Bài báo gốc kiểm chứng SDPO trên ba thiết lập khác nhau. Ở thiết lập RLVR thuần, khi không có sẵn rich feedback (native), họ đã giả lập feedback bằng cách lấy một rollout thành công trong batch làm tham chiếu (reference) cho một rollout thất bại, và SDPO vẫn vượt trội đáng kể so với GRPO, đạt được độ chính xác mà GRPO cần 5 giờ chỉ trong 50 phút trên một tác vụ hóa học {[}1{]}. Ở thiết lập RLRF với execution feedback thực trên bộ LiveCodeBench v6 {[}6{]}, SDPO đạt 48.8\% so với 41.2\% của GRPO, đồng thời bắt kịp độ chính xác cuối cùng của GRPO với tổng số generation ít hơn khoảng 4 lần. Thiết lập thứ ba (test-time discovery trên một bài toán khó duy nhất) chính là phạm vi nghiên cứu của khóa luận này, và sẽ được thảo luận ngay sau đây.

\subsection{Điều gì làm nên SDPO: ablation và scaling}\label{what-drives-sdpo-ablations-and-scaling}

Ba phát hiện từ phần phân tích của bài báo gốc có ảnh hưởng trực tiếp đến các lựa chọn thiết kế trong khóa luận này.

Thứ nhất, cả hai thành phần của SDPO đều thực sự cần thiết. Khi phân rã phương pháp thành các biến thể mức logit, mức token và mức chuỗi, thứ tự hiệu năng logit \textgreater{} token \textgreater{} chuỗi \textgreater{} GRPO {[}1, §4.2{]} cho thấy cả rich feedback và credit assignment dày đều cùng đóng góp vào sự cải thiện, thay vì chỉ một yếu tố đóng vai trò quyết định. Một phân tích cắt bỏ (ablation) theo loại feedback (Bảng 6 của họ) cũng chỉ ra kết luận tương tự: chỉ dùng output môi trường đạt 39.9\% training accuracy, chỉ dùng lời giải trước đó của mô hình đạt 42.6\%, nhưng kết hợp cả hai lại nâng độ chính xác lên tới 48.3\%, nghĩa là hai nguồn tín hiệu này mang tính bổ trợ cho nhau. Ngược lại, nếu đưa attempt thất bại của student vào context của teacher, độ chính xác tụt xuống 44.5\% và entropy giảm 0.23, vì lúc này teacher bị thiên lệch về phía student và mất đi tính đa dạng. Kết quả này chính là một phần động lực cho cách tổ chức teacher-first ở §3.3, vốn chỉ distill các generation của teacher đã qua thanh lọc, thay vì tái sử dụng attempt thất bại của student.

Thứ hai, khả năng tự dạy (self-teaching) là một năng lực đột phá xuất hiện theo quy mô (emergent with scale) {[}1, §4.1{]}. Trên Qwen3-8B, SDPO vượt xa GRPO; trên Qwen3-0.6B hai phương pháp cho kết quả tương đương; còn trên Qwen2.5-1.5B, SDPO lại kém hơn GRPO. Lý do là phương pháp này đòi hỏi năng lực in-context learning đủ mạnh để mô hình được điều kiện hóa thực sự đóng vai trò là một teacher hữu ích. Đặc điểm này đặt mô hình 4B dùng trong khóa luận vào đúng phân khúc mà SDPO được kỳ vọng sẽ phát huy tác dụng, đồng thời gợi ý rằng trên một mô hình mạnh hơn, baseline student-first sẽ tự cải thiện và thu hẹp khoảng cách với teacher-first, một dự đoán sẽ được thảo luận lại ở Chương 6.

Thứ ba, self-teacher được tự mồi (bootstrap) liên tục chứ không cố định ngay từ đầu. Nó được cập nhật đồng thời với student, độ chính xác của nó tăng dần trong quá trình huấn luyện, và cuối cùng student còn vượt qua chính teacher ban đầu; do đó, tiềm năng của phương pháp không bị giới hạn trần bởi teacher khởi điểm {[}1, §4.3{]}. Regularization chính là yếu tố giữ cho quá trình này ổn định: một teacher dùng trust-region (50.6\%) nhỉnh hơn một chút so với teacher dùng EMA (49.3\%) và teacher đóng băng (48.8\%), trong khi một teacher không có regularization sẽ hoàn toàn phân kỳ. SDPO cũng ít bị hiện tượng quên (forgetting) hơn so với việc supervised fine-tuning trực tiếp trên output của self-teacher, xét trên các benchmark held-out (dù có quên nhiều hơn GRPO một chút), cho thấy bản chất on-policy của bước cập nhật giúp bảo toàn năng lực nền tảng tốt hơn {[}1, §4.4{]}. Một biến thể hybrid kết hợp lợi thế của GRPO và SDPO tỏ ra hữu ích cho các mô hình yếu, nhưng không mang lại lợi ích đáng kể cho các mô hình mạnh {[}1, §4.5{]}.

\section{Test-time training và discovery@k}\label{test-time-training-and-discoveryk}

Trong thiết lập test-time {[}1, §5{]}, một bài toán khó duy nhất được cố định, phần thưởng (reward) có dạng nhị phân, và mô hình phải khám phá ra lời giải càng nhanh càng tốt. Thay vì nối một lịch sử ngày càng dài vào context như cách tiếp cận multi-turn reprompting vẫn làm, SDPO nén thẳng feedback vào trọng số bằng cách distill \(\pi_\theta(\cdot \mid x, c)\) vào \(\pi_\theta(\cdot \mid x)\), qua đó tránh được giới hạn của context-window.

Đại lượng đo lường (metric) được sử dụng ở đây là discovery@k {[}1, Def. 5.1{]}: xác suất giải được bài toán trong \(k\) attempt đầu tiên, \(P(T \le k)\) với thời gian khám phá \(T = \min\{k : r(y_k) = 1\}\). Đại lượng này tổng quát hóa pass@k {[}9{]} từ việc lấy mẫu độc lập cùng phân bố (i.i.d.) sang các thuật toán tuần tự có chia sẻ trạng thái hoặc cập nhật trọng số giữa các attempt, và thu về đúng pass@k khi thuật toán chỉ là best-of-k thuần túy. Nó chia sẻ chung nền tảng toán học với hiệu năng time-to-termination của Dolan và Moré {[}12{]}. Về mặt thực nghiệm, SDPO khám phá được lời giải cho 53.2\% các bài toán cực khó, so với 41.5\% của best-of-k, và trên một số bài toán, đây là phương pháp duy nhất có thể giải được, mặc dù độ chính xác của self-teacher ban đầu vẫn dưới 1\% trên gần như mọi bài khó {[}1{]}. Đây chính là điểm cốt lõi nhất: credit assignment dày đặc cho phép mô hình cải thiện dần ngay cả trước khi nó sinh ra được lời giải đúng đầu tiên.

Điều này định vị công trình vào một bức tranh nghiên cứu rộng lớn hơn về inference-time compute. Việc lấy mẫu lặp (repeated sampling) giúp tăng độ phủ (coverage) trên các bài toán khó, mở rộng theo luật lũy thừa (power law) của ngân sách lấy mẫu {[}20{]}; self-consistency tổng hợp nhiều reasoning path đã lấy mẫu thành một đáp án đáng tin cậy hơn {[}18{]}; và test-time scaling có kiểm soát ngân sách sẽ cải thiện khả năng lập luận bằng cách tiêu tốn bổ sung token cho mỗi bài {[}21{]}. Tương đồng nhất về mặt triết lý là phương pháp self-taught reasoning {[}19{]}, tự mồi (bootstrap) một mô hình từ chính các rationale do nó tự sinh ra; phương pháp teacher-first cũng học theo hướng tương tự (từ các quỹ đạo tự sinh đã qua thanh lọc) nhưng distill một target dày đặc theo từng token thay vì fine-tune trực tiếp trên văn bản.

\section{Reprompt template và không gian thiết kế của nó}\label{the-reprompt-template-and-its-design-space}

Hành vi của self-teacher bị chi phối bởi template, thứ quyết định cách \(x\) và \(f\) được định dạng để đưa vào context của nó. Bài báo gốc sử dụng một template dạng slot (Bảng 2 của họ), chèn một rollout thành công trước đó, feedback từ môi trường, và một chỉ dẫn kết thúc; họ ghi nhận rằng hiệu năng "không nhạy với các biến thể cú pháp" của template, trong khi loại feedback lại có ảnh hưởng rõ rệt hơn: một phân tích cắt bỏ (ablation) (Bảng 6 của họ) cho thấy output của môi trường và lời giải của chính mô hình mang tính bổ trợ cho nhau, trong khi việc đưa attempt thất bại của student vào context của teacher lại làm suy giảm tính đa dạng {[}1{]}.

Từ đó mở ra hai khoảng trống nghiên cứu. Bài báo trên chỉ kiểm chứng biến thể cú pháp chứ chưa khảo sát biến thể ngữ nghĩa hay chỉ dẫn, và chính các tác giả cũng nhấn mạnh rằng việc nghiên cứu hệ thống cách reprompt template ảnh hưởng đến hành vi vẫn là hướng mở (future work) {[}1, §7{]}. Khóa luận này tổ chức không gian thiết kế template theo ba chiều: lượng thông tin, thúc đẩy bởi phát hiện của Kim và cộng sự rằng độ giàu của context chi phối mức độ đè nén {[}2{]}; cách đóng khung chỉ dẫn (instruction framing), trực tiếp giải quyết câu hỏi mở vừa nêu; và chiều sâu bộ nhớ (memory depth), thúc đẩy bởi tính tích lũy tuần tự của thiết lập test-time. Bảy template ở §3.4 lấy mẫu đại diện cho các điểm trên ba chiều này, với template mặc định của bài báo gốc đóng vai trò làm điểm neo (anchor).

\section{Epistemic verbalization và sự suy giảm}\label{epistemic-verbalization-and-suppression}

Kim và cộng sự {[}2{]} đưa ra lời cảnh báo đối trọng với SDPO. Họ chỉ ra rằng việc distill từ một context giàu thông tin có thể gây phản tác dụng: nó làm suy giảm năng lực của mô hình bằng cách kìm hãm các biểu đạt nhận thức (epistemic verbalization) của chính nó, tức các dấu hiệu chỉ sự bất định (uncertainty markers) như "wait" hay "maybe" vốn thường đi kèm với lập luận chain-of-thought {[}17{]}. Phân tích của họ liên kết biên độ của sự suy giảm này với lượng thông tin tương hỗ \(I(y; c \mid x)\) giữa output và context: context càng giàu thông tin, mức độ suy giảm càng mạnh. Họ khảo sát phổ ảnh hưởng này bằng cách thay đổi nội dung context của teacher, từ rỗng (\(c = \varnothing\)), qua một lời giải đã được lược bỏ thinking trace kèm một đáp án được sinh lại, cho tới một lời giải đầy đủ, và nhận thấy mức độ suy giảm tỉ lệ thuận với lượng thông tin mà context mang lại. Hiệu ứng này còn tích lũy qua các bước distillation nối tiếp nhau, nên họ khuyến nghị đóng băng teacher (không cập nhật EMA) để phá vỡ vòng lặp feedback, một khuyến nghị đi ngược lại định hướng của bài báo gốc, vốn ưu tiên một teacher được cập nhật liên tục và được điều chuẩn (regularized) (§2.2.4). Nhìn rộng hơn, các mô hình ngôn ngữ vốn đã được hiệu chuẩn (calibrated) một phần về giới hạn tri thức của chúng {[}23{]}, nên việc triệt tiêu các marker biểu thị sự bất định này là một con đường khá hợp lý dẫn tới sự suy giảm năng lực.

Đối với khóa luận này, sự liên quan nằm ở hai điểm. Một, rủi ro suy giảm biểu đạt trở nên rõ rệt nhất ngay tại thời điểm context chứa sẵn đáp án, chính là math leak regime được nghiên cứu trong pilot toán (§4.9), nơi teacher hoàn toàn có thể sao chép đáp án thay vì thực sự suy dẫn ra nó. Hai, tính chất tích lũy qua các bước chính là điều mà thiết lập test-time thực thi lặp đi lặp lại, vì mỗi bước TTT đều distill từ cùng một context mang-đáp-án (answer-bearing context).

\section{SDFT và mối lo về rò rỉ (leak)}\label{sdft-and-the-leak-concern}

Shenfeld và cộng sự {[}3{]} đề xuất self-distillation fine-tuning (SDFT) cho continual learning, một biến thể distillation on-policy nhằm tối thiểu hóa sự thay đổi (minimum-change), hướng theo các generation được điều kiện hóa của chính mô hình. Đây là một phương pháp tương đồng với SDPO: cũng dùng mô hình làm teacher của chính nó, nhưng mục đích lại là bảo toàn năng lực nền tảng trong khi thích nghi với tri thức mới, do đó việc giảm thiểu các thay đổi ngoài ý muốn trở thành một mục tiêu thiết kế rõ ràng. Cách tổ chức teacher-first của khóa luận này có nét tương đồng với SDFT ở chỗ nó chưng cất (distill) hướng theo các quỹ đạo tự sinh đã qua thanh lọc (filtered model-generated trajectories), nhưng lại sử dụng execution feedback làm context điều kiện hóa thay vì một demonstration cố định (§3.3.1). Khuôn khổ tối thiểu hóa thay đổi của SDFT cũng làm rõ nét hơn câu hỏi về rò rỉ (leak): nếu distillation truyền tải nhiều thông tin hơn mức dự kiến, thì thứ thực sự đang bị sao chép là gì, đây chính là câu hỏi mà pilot toán sẽ trả lời một cách rất cụ thể.

\section{Khoảng trống nghiên cứu mà khóa luận giải quyết}\label{the-gap-this-thesis-addresses}

Bài báo gốc nghiên cứu thuật toán SDPO on-policy, student-first, và tập trung phần lớn phân tích vào các chế độ train-time; công trình này có đề cập đến thiết lập test-time nhưng chưa khảo sát reprompt template trong đó, cũng như chưa đo lường hành vi epistemic. Kim và cộng sự đã mô tả đặc trưng hiện tượng suy giảm biểu đạt nhận thức (epistemic suppression), nhưng không đặt trong bối cảnh khám phá mã (code-discovery) tại thời điểm suy luận. Trong khi đó, Shenfeld và cộng sự tối ưu hóa nhằm giảm thiểu sự thay đổi cho continual learning, thay vì hướng đến mục tiêu tối đa hóa khả năng khám phá lời giải (discovery) trên một bài toán khó duy nhất.

Do đó, khoảng trống nghiên cứu bao gồm ba khía cạnh. Thứ nhất là việc đề xuất một cách tổ chức teacher-first cho phương pháp execution-guided self-distillation, chưng cất (distill) từ các generation của teacher đã qua thanh lọc, đặt phương pháp ở vị trí giao thoa giữa SDPO và SFT, nhằm trả lời câu hỏi liệu nó có thực sự tăng tốc test-time discovery hay không (RO-method). Thứ hai là tác động của cách định dạng (formulate) reprompt-template đối với quá trình discovery, đây chính là câu hỏi mở mà bài báo gốc đã chỉ ra (RO1). Thứ ba là việc đặc trưng hóa nhằm xác định khi nào cách tiếp cận này phát huy hiệu quả, được đóng khung dưới dạng ranh giới quy trình-và-giá trị (value-versus-procedure) mà sự tương phản giữa code-và-toán vạch rõ ra. Phần còn lại của khóa luận sẽ lần lượt giải quyết ba khoảng trống nghiên cứu này.
